% \documentclass[aspectratio=169,notes]{beamer}
\documentclass[aspectratio=169]{beamer}
\usetheme[faculty=phil]{fibeamer}
\usepackage{polyglossia}
\setmainlanguage{english} %% main locale instead of `english`, you
%% can typeset the presentation in either Czech or Slovak,
%% respectively.
\setotherlanguages{russian} %% The additional keys allow
%%
%%   \begin{otherlanguage}{czech}   ... \end{otherlanguage}
%%   \begin{otherlanguage}{slovak}  ... \end{otherlanguage}
%%
%% These macros specify information about the presentation
\title[Introduction to Robotics (In2Ro)]{Introduction to Robotics, Lecture 4} %% that will be typeset on the
\subtitle{Differential Kinematics 
\\ Joint and Tool Velocities \\ \  } %% title page.
\author{Oleg Bulichev}
%% These additional packages are used within the document:
\usepackage{ragged2e}  % `\justifying` text
\usepackage{booktabs}  % Tables
\usepackage{tabularx}
\usepackage{tikz}      % Diagrams
\usetikzlibrary{calc, shapes, backgrounds, trees, positioning, arrows.meta}
\usetikzlibrary{decorations.pathreplacing,calligraphy,calc,graphs}
\usepackage{amsmath, amssymb}
\usepackage{url}       % `\url`s
\usepackage{listings}  % Code listings
\usepackage{floatrow}
\usepackage{subcaption}
\usepackage{mathtools}
\usepackage{todonotes}
\usepackage{fontspec}
\usepackage{multicol}
\usepackage{pdfpages}
\usepackage{wrapfig}
\usepackage{qrcode}
\usepackage{animate}
\usepackage{booktabs}
\usepackage{multirow}
\usepackage{graphicx}
\usepackage{colortbl}

\graphicspath{{resources/}}
\frenchspacing

\setbeamertemplate{caption}{\raggedright\insertcaption\par}
\usetikzlibrary{graphs}

% \usepackage[backend=biber,style=ieee,autocite=footnote]{biblatex}
% \addbibresource{biblio.bib}
% \DefineBibliographyStrings{english}{%
%   bibliography = {References},}

\newcommand{\oleg}[2][] {\todo[color=red, #1] {OLEG:\\ #2}}
\newcommand{\fbckg}[1]{\usebackgroundtemplate{\includegraphics[width=\paperwidth]{#1}}}%frame background

\usepackage[framemethod=TikZ]{mdframed}
\newcommand{\dbox}[1]{
\begin{mdframed}[roundcorner=3pt, backgroundcolor=yellow, linewidth=0]
\vspace{1mm}
{#1}
\vspace{1mm}
\end{mdframed}
}

\begin{document}
\setlength{\abovedisplayskip}{0pt}
\setlength{\belowdisplayskip}{0pt}
\setlength{\abovedisplayshortskip}{0pt}
\setlength{\belowdisplayshortskip}{0pt}

\fbckg{fibeamer/figs/title_page.png}
\frame[c]{\setcounter{framenumber}{0}
    \usebeamerfont{title}%
    \usebeamercolor[fg]{title}%
    \begin{minipage}[b][6.5\baselineskip][b]{\textwidth}%
        \textcolor{black}{\raggedright\inserttitle}
    \end{minipage}
    % \vskip-1.5\baselineskip

    \usebeamerfont{subtitle}%
    \usebeamercolor[fg]{framesubtitle}%
    \begin{minipage}[b][3\baselineskip][b]{\textwidth}
        \raggedright%
        \insertsubtitle%
    \end{minipage}
    \vskip.25\baselineskip
}

\fbckg{fibeamer/figs/common.png}

\begin{frame}[t]{Differential Kinematics}
    Differential kinematics relates joint velocities to the linear and
    angular velocities of a robot's end effector. This relation is used
    to control motion, construct trajectories, and analyze singular
    configurations.
    \bigskip
    The following three sources explain the topic in different ways.
    Start with the clearest explanation for you, then compare all three.
\end{frame}

\begin{frame}[t]{Course Material Robotics 1}
    Prof. Alessandro De Luca — Sapienza University of Rome.
    \medskip
    \begin{itemize}
        \item \href{https://www.diag.uniroma1.it/deluca/rob1_en/11_DifferentialKinematics.pdf}{11. Differential Kinematics — PDF}
    \end{itemize}
\end{frame}

\begin{frame}[t]{Introduction to Robotics — Stanford}
    Prof. Oussama Khatib — Stanford Engineering Everywhere.
    \medskip
    \begin{itemize}
        \item \href{https://www.youtube.com/watch?v=fwHc0a8DMQ0&list=PL65CC0384A1798ADF&index=6}{Lecture 6: The Jacobian and instantaneous kinematics}
        \item \href{https://www.youtube.com/watch?v=6SRTAoyzC6A&list=PL65CC0384A1798ADF&index=7}{Lecture 7: Jacobians and singularities}
    \end{itemize}
\end{frame}

\begin{frame}[t]{Modern Robotics — Chapter 5}
    Kevin M. Lynch and Frank C. Park — Northwestern University.
    \medskip
    \begin{itemize}
        \item \href{https://www.youtube.com/watch?v=KbI8HN3imtQ&list=PLggLP4f-rq02vX0OQQ5vrCxbJrzamYDfx&index=26}{Chapter 5.1.1: Space Jacobian}
        \item \href{https://www.youtube.com/watch?v=69HeL2XvL9Y&list=PLggLP4f-rq02vX0OQQ5vrCxbJrzamYDfx&index=27}{Chapter 5.1.2: Body Jacobian}
    \end{itemize}
\end{frame}

\begin{frame}[t]{Reference Materials}
    \begin{itemize}
        \item A. De Luca --- 11 DifferentialKinematics
        \item Oussama Khatib --- Differential kinematics
        \item Kevin M. Lynch and Frank C. Park --- Space and body Jacobians
    \end{itemize}
\end{frame}


\fbckg{fibeamer/figs/last_page.png}
\frame[plain]{}

\end{document}
